\documentclass[11pt,a4paper]{article}
\usepackage[utf8]{inputenc}
\usepackage[T1]{fontenc}
\usepackage[spanish]{babel}
\usepackage{geometry}
\usepackage{amsmath}
\usepackage{amssymb}
\usepackage{graphicx}
\usepackage{enumitem}
\usepackage{booktabs}
\usepackage{array}
\usepackage{multirow}
\usepackage{fancyhdr}
\usepackage{xcolor}
\usepackage{mdframed}
\usepackage{tikz}
\usepackage{pgfplots}
\pgfplotsset{compat=1.18}

\geometry{top=2cm, bottom=2.5cm, left=2.5cm, right=2.5cm}

\pagestyle{fancy}
\fancyhf{}
\rhead{Tecnolog\'ia Digital VI --- 1er Semestre 2026}
\lhead{Parcial --- Modelo 1}
\rfoot{P\'agina \thepage}

\definecolor{tdblue}{RGB}{0,74,135}
\definecolor{lightgray}{RGB}{245,245,245}

\newmdenv[linecolor=tdblue, linewidth=1.5pt, roundcorner=4pt, backgroundcolor=lightgray, innertopmargin=6pt, innerbottommargin=6pt]{infobox}

\newcommand{\puntaje}[1]{\textbf{[#1 pts]}}

\begin{document}

%--- Encabezado ---
\begin{center}
  {\Large\bfseries Universidad Torcuato Di Tella}\\[4pt]
  {\large Tecnolog\'ia Digital VI --- Parcial Modelo 1}\\[2pt]
  {\normalsize 1er Semestre 2026}\\[8pt]
\end{center}

\begin{infobox}
\textbf{Instrucciones:} El parcial tiene una duraci\'on de \textbf{90 minutos}. Est\'a prohibido el uso de computadoras o c\'odigo. S\'i se permite el uso de una hoja de f\'ormulas. Cada secci\'on indica el puntaje. Puntaje total: \textbf{100 puntos}.
\end{infobox}

\vspace{0.5cm}

Nombre y apellido: \underline{\hspace{8cm}} \hfill Legajo: \underline{\hspace{3cm}}

\vspace{0.8cm}

%=============================================================
\section*{Secci\'on 1: Verdadero o Falso \puntaje{24}}
%=============================================================

\noindent Indique \textbf{V} (Verdadero) o \textbf{F} (Falso) para cada afirmaci\'on. Si la afirmaci\'on es falsa, \textbf{justifique brevemente} por qu\'e (sin justificaci\'on no se otorga puntaje en los falsos).

\vspace{0.3cm}

\begin{enumerate}[label=\textbf{\arabic*.}, leftmargin=*, itemsep=18pt]

\item \underline{\hspace{1cm}} En un mecanismo MCAR (\textit{Missing Completely at Random}), la ausencia de un valor depende del valor faltante en s\'i mismo, lo que lo hace el tipo de dato faltante m\'as problem\'atico.

\vspace{0.5cm}
\noindent\textit{Justificaci\'on (si es F):} \underline{\hspace{12cm}}

\vspace{0.2cm}
\underline{\hspace{14cm}}

\item \underline{\hspace{1cm}} Aplicar un \textit{StandardScaler} ajustado sobre el conjunto de \textit{test} antes de dividir los datos en train/test constituye un caso de \textit{data leakage}.

\vspace{0.5cm}
\noindent\textit{Justificaci\'on (si es F):} \underline{\hspace{12cm}}

\vspace{0.2cm}
\underline{\hspace{14cm}}

\item \underline{\hspace{1cm}} En K-Fold Cross-Validation con $k=5$, cada muestra del dataset es usada exactamente una vez como muestra de validaci\'on a lo largo de las 5 iteraciones.

\vspace{0.5cm}
\noindent\textit{Justificaci\'on (si es F):} \underline{\hspace{12cm}}

\vspace{0.2cm}
\underline{\hspace{14cm}}

\item \underline{\hspace{1cm}} La Regresi\'on Log\'istica es un modelo no lineal porque utiliza la funci\'on sigmoide para transformar su salida.

\vspace{0.5cm}
\noindent\textit{Justificaci\'on (si es F):} \underline{\hspace{12cm}}

\vspace{0.2cm}
\underline{\hspace{14cm}}

\item \underline{\hspace{1cm}} La regularizaci\'on L1 (Lasso) tiende a producir coeficientes exactamente iguales a cero, lo que la convierte en un m\'etodo impl\'icito de selecci\'on de variables.

\vspace{0.5cm}
\noindent\textit{Justificaci\'on (si es F):} \underline{\hspace{12cm}}

\vspace{0.2cm}
\underline{\hspace{14cm}}

\item \underline{\hspace{1cm}} Un \'arbol de decisi\'on muy profundo (sin regularizaci\'on) suele tener bajo sesgo y baja varianza, lo que lo convierte en un modelo ideal sin necesidad de ajuste.

\vspace{0.5cm}
\noindent\textit{Justificaci\'on (si es F):} \underline{\hspace{12cm}}

\vspace{0.2cm}
\underline{\hspace{14cm}}

\item \underline{\hspace{1cm}} Un AUC-ROC de 0.5 indica que el modelo tiene un rendimiento equivalente al de clasificar al azar.

\vspace{0.5cm}
\noindent\textit{Justificaci\'on (si es F):} \underline{\hspace{12cm}}

\vspace{0.2cm}
\underline{\hspace{14cm}}

\item \underline{\hspace{1cm}} El \'indice Gini y la Entrop\'ia, como funciones de impureza en \'arboles de decisi\'on, alcanzan su valor m\'aximo cuando todas las muestras de un nodo pertenecen a la misma clase.

\vspace{0.5cm}
\noindent\textit{Justificaci\'on (si es F):} \underline{\hspace{12cm}}

\vspace{0.2cm}
\underline{\hspace{14cm}}

\end{enumerate}

\vspace{0.5cm}

%=============================================================
\section*{Secci\'on 2: M\'ultiple Choice \puntaje{24}}
%=============================================================

\noindent Marque con un c\'irculo la \textbf{\'unica} opci\'on correcta en cada pregunta.

\vspace{0.3cm}

\begin{enumerate}[label=\textbf{\arabic*.}, leftmargin=*, itemsep=20pt]

\item Un dataset de detecci\'on de fraude tiene 98\% de transacciones legales y 2\% de transacciones fraudulentas. Un modelo que siempre predice ``legal'' obtiene un \textit{Accuracy} del 98\%. ¿Cu\'al de las siguientes afirmaciones describe mejor la situaci\'on?

\begin{enumerate}[label=\alph*)]
  \item El modelo es excelente; un 98\% de accuracy es muy alto para cualquier problema.
  \item El accuracy es enga\~noso en este caso; m\'etricas como F1-Score o Precision-Recall son m\'as adecuadas para evaluar el rendimiento real.
  \item El problema se resolver\'ia aumentando el umbral de decisi\'on del modelo a 0.99.
  \item La soluci\'on correcta es aplicar un StandardScaler antes de entrenar.
\end{enumerate}

\item ¿Cu\'al es el efecto principal de aumentar el par\'ametro \texttt{max\_depth} en un \'arbol de decisi\'on?

\begin{enumerate}[label=\alph*)]
  \item Reduce el tiempo de entrenamiento y aumenta el sesgo del modelo.
  \item Aumenta la complejidad del modelo, reduciendo el sesgo pero potencialmente aumentando la varianza (overfitting).
  \item Reduce tanto el sesgo como la varianza, mejorando siempre el desempe\~no.
  \item No tiene efecto sobre el overfitting; solo afecta la velocidad de predicci\'on.
\end{enumerate}

\item Un equipo aplica \textit{Label Encoding} a la variable ``Pa\'is'' con valores \{Argentina=0, Brasil=1, Chile=2\} y luego entrena una regresi\'on log\'istica. ¿Cu\'al es el problema principal de este enfoque?

\begin{enumerate}[label=\alph*)]
  \item El modelo no puede manejar m\'as de dos categor\'ias en una variable.
  \item El encoding impone un orden num\'erico artificial (Argentina < Brasil < Chile) que el modelo interpretar\'a como una relaci\'on ordinal inexistente.
  \item El Label Encoding solo funciona con \'arboles de decisi\'on, no con modelos lineales.
  \item No hay ning\'un problema; Label Encoding es siempre equivalente a One-Hot Encoding.
\end{enumerate}

\item La curva Precision-Recall es preferida sobre la curva ROC cuando:

\begin{enumerate}[label=\alph*)]
  \item El dataset est\'a perfectamente balanceado entre clases.
  \item Se trabaja con un problema de regresi\'on, no de clasificaci\'on.
  \item Las clases est\'an muy desbalanceadas y la clase positiva es la de inter\'es, ya que la curva ROC puede resultar optimista al verse influida por la gran cantidad de verdaderos negativos.
  \item El modelo utilizado es un \'arbol de decisi\'on en lugar de una regresi\'on log\'istica.
\end{enumerate}

\item ¿Qu\'e representa el error \textit{Out-of-Bag} (OOB) en Random Forest?

\begin{enumerate}[label=\alph*)]
  \item El error cometido sobre el conjunto de test reservado al final del proceso.
  \item Una estimaci\'on del error de generalizaci\'on calculada usando, para cada muestra, solo los \'arboles que \textbf{no} la incluyeron en su muestra de bootstrap.
  \item El error promedio de entrenamiento de todos los \'arboles individuales.
  \item El error producido por usar demasiados \'arboles, causando overfitting en el ensamble.
\end{enumerate}

\item ¿Cu\'al de las siguientes opciones describe correctamente la diferencia entre \textit{Bagging} y \textit{Boosting}?

\begin{enumerate}[label=\alph*)]
  \item Bagging entrena modelos secuencialmente ponderando los errores; Boosting los entrena en paralelo sobre distintas muestras bootstrap.
  \item Bagging entrena modelos en paralelo sobre muestras bootstrap independientes para reducir la varianza; Boosting entrena modelos secuencialmente, donde cada modelo se focaliza en los errores del anterior, reduciendo el sesgo.
  \item Bagging y Boosting son t\'ecnicamente id\'enticos, solo difieren en el lenguaje de programaci\'on utilizado.
  \item Bagging solo funciona con modelos lineales; Boosting solo funciona con \'arboles.
\end{enumerate}

\end{enumerate}

\vspace{0.5cm}

%=============================================================
\section*{Secci\'on 3: Interpretaci\'on de Gr\'aficos \puntaje{20}}
%=============================================================

\subsection*{Ejercicio A: Curvas de aprendizaje \puntaje{10}}

El siguiente gr\'afico muestra las curvas de error de entrenamiento y validaci\'on de un modelo en funci\'on de la profundidad m\'axima del \'arbol (\texttt{max\_depth}).

\vspace{0.3cm}
\begin{center}
\begin{tikzpicture}
\begin{axis}[
  width=10cm, height=6cm,
  xlabel={Profundidad m\'axima (\texttt{max\_depth})},
  ylabel={Error},
  xmin=1, xmax=10,
  ymin=0, ymax=1,
  xtick={1,2,3,4,5,6,7,8,9,10},
  legend pos=north east,
  grid=major,
  grid style={line width=0.3pt, draw=gray!30}
]
\addplot[blue, thick, mark=o] coordinates {
  (1,0.80)(2,0.65)(3,0.50)(4,0.38)(5,0.28)(6,0.20)(7,0.14)(8,0.10)(9,0.07)(10,0.05)
};
\addlegendentry{Error de Entrenamiento}
\addplot[red, thick, mark=square] coordinates {
  (1,0.82)(2,0.66)(3,0.52)(4,0.42)(5,0.38)(6,0.40)(7,0.45)(8,0.52)(9,0.58)(10,0.65)
};
\addlegendentry{Error de Validaci\'on}
\end{axis}
\end{tikzpicture}
\end{center}

\vspace{0.3cm}

\begin{enumerate}[label=\textbf{\alph*)}]
  \item ¿En qu\'e valor de \texttt{max\_depth} el modelo comienza a presentar \textit{overfitting}? Justifique observando el comportamiento de ambas curvas.

  \vspace{2.5cm}

  \item ¿Qu\'e ocurre con el modelo cuando \texttt{max\_depth} = 1 o 2? ¿C\'omo se llama este fen\'omeno?

  \vspace{2.5cm}

  \item Si deber\'ia elegir un valor de \texttt{max\_depth} para producir el mejor modelo, ¿cu\'al elegiria y por qu\'e?

  \vspace{2.5cm}
\end{enumerate}

\subsection*{Ejercicio B: Matriz de confusi\'on \puntaje{10}}

Un modelo de clasificaci\'on binaria para detectar correos spam (positivo = spam) produjo la siguiente matriz de confusi\'on sobre el conjunto de test:

\vspace{0.3cm}
\begin{center}
\begin{tabular}{cc|cc}
  & & \multicolumn{2}{c}{\textbf{Predicci\'on}} \\
  & & Spam & No Spam \\
  \hline
  \multirow{2}{*}{\textbf{Real}} & Spam    & 80  & 20  \\
                                  & No Spam & 30  & 870 \\
\end{tabular}
\end{center}

\vspace{0.3cm}

\begin{enumerate}[label=\textbf{\alph*)}]
  \item Calcule el \textbf{Recall} (sensibilidad) para la clase Spam y explique con palabras qu\'e nos indica este valor en el contexto del problema.

  \vspace{2.5cm}

  \item Calcule la \textbf{Precision} para la clase Spam e interprete el resultado.

  \vspace{2.5cm}

  \item En este contexto (filtro de spam), ¿es m\'as grave tener un Recall bajo o una Precision baja? Justifique considerando las consecuencias pr\'acticas de cada tipo de error.

  \vspace{2.5cm}
\end{enumerate}

\vspace{0.3cm}

%=============================================================
\section*{Secci\'on 4: Desarrollo \puntaje{32}}
%=============================================================

\begin{enumerate}[label=\textbf{\arabic*.}, leftmargin=*, itemsep=10pt]

\item \puntaje{10} \textbf{Preprocesamiento y Data Leakage}

Un equipo de Data Science tiene un pipeline para predecir el precio de propiedades. El pipeline que implementaron es el siguiente:

\begin{enumerate}[label=\Roman*.]
  \item Cargar el dataset completo (10.000 propiedades).
  \item Imputar valores faltantes con la \textit{media} de cada columna, calculada sobre todo el dataset.
  \item Escalar todas las variables num\'ericas con MinMaxScaler ajustado sobre todo el dataset.
  \item Dividir en 80\% train y 20\% test aleatoriamente.
  \item Entrenar un modelo de Regresi\'on Lineal sobre el conjunto de train.
  \item Evaluar con RMSE sobre el conjunto de test.
\end{enumerate}

\begin{enumerate}[label=\textbf{\alph*)}]
  \item Identifique \textbf{todos} los errores conceptuales presentes en el pipeline. Explique por qu\'e cada uno es un problema.
  \vspace{4cm}
  \item Describa el pipeline corregido, paso a paso, que elimina los errores identificados.
  \vspace{4cm}
\end{enumerate}

\item \puntaje{10} \textbf{\'Arboles de Decisi\'on y Regularizaci\'on}

\begin{enumerate}[label=\textbf{\alph*)}]
  \item Explique intuitivamente por qu\'e los \'arboles de decisi\'on son modelos con \textbf{alta varianza}. ¿Qu\'e propiedad del algoritmo de construcci\'on los hace tan sensibles a cambios peque\~nos en los datos?
  \vspace{3.5cm}
  \item Mencione \textbf{tres} hiperpar\'ametros que se pueden ajustar para regularizar un \'arbol de decisi\'on (reducir overfitting) y explique brevemente el efecto de cada uno.
  \vspace{3.5cm}
  \item Un \'arbol de regresi\'on fue entrenado con valores de temperatura entre 10\textdegree C y 40\textdegree C. Hoy en producci\'on recibe una observaci\'on con temperatura = 55\textdegree C. ¿Qu\'e valor predecir\'a el modelo? ¿Es esto un problema? ¿Por qu\'e?
  \vspace{3cm}
\end{enumerate}

\item \puntaje{12} \textbf{Caso integral: Sistema de recomendaci\'on de cr\'editos}

Un banco quiere construir un sistema para predecir si un cliente pagar\'a o no un cr\'edito (variable objetivo: \texttt{default} = 1 si no paga, 0 si paga). El dataset tiene:
\begin{itemize}
  \item 500.000 clientes hist\'oricos, de los cuales solo 12.000 no pagaron.
  \item Variables num\'ericas: ingreso mensual, edad, monto del cr\'edito, historial de pagos.
  \item Variables categ\'oricas: tipo de empleo, zona geogr\'afica, nivel de educaci\'on.
  \item Un 8\% de valores nulos en ``ingreso mensual'' y un 3\% en ``historial de pagos''.
\end{itemize}

El equipo propone entrenar directamente un \'arbol de decisi\'on con par\'ametros por defecto y reportar el Accuracy al directorio.

\begin{enumerate}[label=\textbf{\alph*)}]
  \item ¿Qu\'e problemas tiene el dataset que deben resolverse antes de entrenar? Mencione al menos \textbf{tres} y proponga una estrategia para cada uno.
  \vspace{4cm}
  \item ¿Por qu\'e reportar el Accuracy al directorio es enga\~noso en este caso? ¿Qu\'e m\'etrica(s) recomendar\'ia y por qu\'e?
  \vspace{3.5cm}
  \item El equipo propone un \'arbol de decisi\'on. ¿Qu\'e argumentos dar\'ia para considerar tambi\'en un \textbf{Random Forest} o \textbf{Gradient Boosting}? ¿Existe alguna raz\'on para mantener el \'arbol simple?
  \vspace{3.5cm}
\end{enumerate}

\end{enumerate}

\vspace{1cm}
\begin{center}
\textit{--- Fin del Parcial Modelo 1 ---}
\end{center}

\end{document}
